\chapter{Methodology}
\label{cha:methodology}

This chapter presents the methodological framework developed to generate functionality-preserving adversarial traffic capable of evading modern NIDS. The proposed approach extends an existing GAN–PSO-based adversarial traffic manipulator by introducing \textbf{FRANTIC} (FRAgmeNTation adversarIal Component), a stochastic structural transformation layer operating at the IPv4 fragmentation level.

\section{Motivations}
The original manipulator, as discussed in Section~\ref{sec:existing_manipulator}, represents a foundational approach to adversarial traffic generation in raw packet space. By combining GAN-guided boundary exploration with PSO-based mutation operators, it enables realistic evasion of ML-based NIDS while maintaining protocol validity and attack semantics. However, the analysis in Section~\ref{sec:limitations_existing} revealed a fundamental limitation: the framework operates predominantly in feature space, manipulating header-level and temporal characteristics (e.g., inter-arrival times), while leaving the structural aspect of traffic largely untouched.

In particular, transformations at the IP layer, such as packet fragmentation, are not considered. Yet fragmentation significantly influences packet-level and flow-level statistical descriptors, burst structures, and fixed-size tensor representations commonly exploited by modern ML-NIDS. Despite its historical relevance in intrusion evasion, this structural dimension remains absent from the GAN–PSO design, thereby leaving a substantial portion of the adversarial search space unexplored.

To address this gap, FRANTIC is introduced as a post-manipulation transformation stage. Rather than replacing the existing optimization strategy, it augments it by enabling adversarial exploration along two complementary axes:

\begin{itemize}
    \item \textbf{Temporal/Statistical Axis}: Optimization of inter-arrival times and payload-level perturbations via GAN-guided PSO.
    \item \textbf{Structural Axis}: Stochastic IPv4-layer fragmentation to alter packet size distributions, flow geometry, and structural entropy.
\end{itemize}

By jointly operating across these axes, the resulting framework enables wider adversarial exploration in feature space. This significantly broadens the reachable adversarial manifold while preserving strict protocol correctness and end-to-end attack functionality. The following sections detail the architectural design, formalization, and implementation of the proposed fragmentation-aware manipulator.

\section{System Architecture}
\label{sec:system_overview}
The proposed adversarial generation pipeline operates in two distinct and sequential phases: (i) a feature-space manipulation phase, and (ii) a structural fragmentation phase.

\begin{figure}[h]
    \centering
    \includegraphics[width=\linewidth]{IMG/FlowChart.png}
    \caption{Overview of the adversarial dataset generation pipeline.}
    \label{fig:pipeline}
\end{figure}

As illustrated in Figure~\ref{fig:pipeline}, the original CIC-IDS2019 dataset is separated into benign and malicious traffic. Malicious flows undergo traffic manipulation followed by stochastic IP fragmentation, while benign traffic remains unchanged. The processed malicious samples are then recombined with benign traffic to construct the final adversarial dataset.
\\\\\\\\
\subsection{Phase 1: Feature-Space Optimization (GAN + PSO)}

This phase implements the original manipulator logic:

\begin{itemize}
    \item A GAN is trained to approximate the benign feature distribution.
    \item A PSO algorithm iteratively perturbs malicious flows, primarily via temporal adjustments and payload-level modifications, minimizing detection confidence.
\end{itemize}

The output consists of optimized malicious PCAPs that remain structurally unchanged at the IP layer but exhibit adversarial statistical properties.

\subsection{Phase 2: Stochastic Fragmentation}

The fragmentation stage applies probabilistic, controlled, RFC-compliant IPv4 fragmentation as a post-processing transformation. By decomposing packets into randomized yet protocol-valid fragments, the system induces measurable shifts in packet- and flow-level statistical features while preserving end-to-end attack functionality.

This structural perturbation challenges flow-based ML-NIDS that rely on packet size distributions, burst characteristics, and fixed-length tensor representations, thereby expanding the adversarial traffic space without compromising protocol correctness.

\section{FRANTIC}
\label{sec:fragmentation_module}
The fragmentation module addresses the structural rigidity of adversarial samples generated by purely feature-based manipulation.

FRANTIC operates at the IPv4 layer on packets generated by the manipulation stage. Fragmentation is applied selectively and under strict protocol constraints to ensure semantic transparency and transport-layer stability.

Network traffic features are abstractions derived from packet streams. By altering stream structure via fragmentation, we effectively perturb the mapping:

\[
\text{Traffic Space} \longrightarrow \text{Feature Space}.
\]

\subsection{Design Principles}

The module is governed by four core design principles:

\begin{enumerate}
    \item \textbf{Strict RFC Compliance}: All fragments conform to RFC~791, including 8-byte alignment, correct flag handling, and consistent IP Identification values.
    \item \textbf{Configurability}: Parameters controlling fragmentation probability and fragment size bounds enable systematic experimentation.
    \item \textbf{Semantic Transparency}: Fragmentation must remain invisible to transport and application layers; reassembly at the destination must fully restore the malicious payload.
    \item \textbf{Stochastic Randomization}: Randomized selection and fragment sizes prevent deterministic evasion signatures.
\end{enumerate}

\subsection{Architecture}
FRANTIC is divided into two logical layers: the Orchestration Layer and the Reconstruction Layer.

The first phase is implemented in \texttt{manipulator.py}. This component iterates over the packet stream, maintains global configuration parameters, and determines packet eligibility for subsequent processing. The complete source code is provided in Appendix~\ref{sec:manipulator_code}.

The configuration parameters include:

\begin{itemize}
    \item $\rho$ (\texttt{fragment\_prob}): Fragmentation probability, $0 \le \rho \le 1$.
    \item $S_{min}$: Minimum fragment size (default: 64 bytes).
    \item $S_{max}$: Maximum fragment size (default: 500 bytes).
\end{itemize}

The reconstruction layer is implemented in \texttt{rebuilder.py} (Appendix~\ref{sec:rebuilder_code}). This component performs byte-level payload decomposition and header reconstruction, ensuring syntactic correctness, accurate fragment offset computation, preservation of protocol metadata, and full reassembly consistency at the destination.

\subsection{Algorithmic Formulation}
\label{sec:algorithm}

The procedure operates on a packet stream $P$ and produces a transformed stream $P'$ in which eligible IPv4 packets may be decomposed into multiple protocol-valid fragments while preserving semantic correctness.
\\

Let the original packet stream be:

\[
P = \{p_1, p_2, \dots, p_n\}.
\]

After fragmentation, the resulting stream becomes:

\[
P' = \{f_{1,1}, f_{1,2}, \dots, f_{n,k}\}.
\]

where each packet $p_i$ may correspond to one or multiple fragments $f_{i,j}$.
\\

Packet eligibility requires:

\[
L_{payload} > S_{min}
\quad \text{and} \quad
r \sim \mathcal{U}[0,1], \quad r \le \rho.
\]

Only IPv4 packets whose payload exceeds a configurable minimum threshold are considered. Small control packets, such as ACK-only TCP segments, are excluded in order to preserve connection stability and avoid unintended side effects on transport-layer behavior.
Let $r$ be a realization of a continuous uniform random variable $r \sim \mathcal{U}(0,1)$. Fragmentation is applied if $r \le \rho$, thereby implementing Bernoulli selection with parameter $\rho$.
Only packets satisfying all conditions enter the fragmentation routine.
\\

For eligible packets, fragmentation is applied probabilistically with $P_{frag} \in [0,1]$, iteratively over the remaining payload length $L_{rem}$. This stochastic selection prevents deterministic transformation patterns and introduces variability across executions, thereby increasing structural entropy within the traffic stream.

When a packet is selected, fragment sizes are sampled from a bounded interval $[S_{min}, S_{max}]$. Intermediate fragments are aligned to 8-byte boundaries to comply with IPv4 offset encoding requirements, while the final fragment absorbs the remaining payload. 
\\

Intermediate fragment size is sampled as:

\[
s_{frag} \sim \mathcal{U}(S_{min}, s_{limit}),
\quad
s_{limit} = \min(S_{max}, L_{rem} - S_{min}).
\]

This ensures that a valid final fragment of at least $S_{min}$ bytes can still be generated.
\\

For non-final fragments:

\[
s_{frag} \leftarrow \left\lfloor \frac{s_{frag}}{8} \right\rfloor \times 8.
\]
\\\\
To formalize the structural transformation process, Algorithm~\ref{alg:stochastic_fragmentation} presents the stochastic RFC-compliant fragmentation procedure. The algorithm abstracts from implementation-specific details and describes the high-level control flow governing packet eligibility, probabilistic selection, fragment size generation, and RFC-compliant header construction.

\begin{algorithm}[h]
\caption{Stochastic RFC-Compliant IP Fragmentation}
\label{alg:stochastic_fragmentation}
\begin{algorithmic}[1]
\Require Packet stream $P$, probability $\rho$, $S_{min}$, $S_{max}$
\Ensure Fragmented packet stream $P'$
\State $P' \gets \emptyset$
\For{each packet $p \in P$}
    \If{$p$ is IPv4 \textbf{and} payload$(p) > S_{min}$ \textbf{and} rand$(0,1) \le \rho$}
        \State $ID \gets$ random 16-bit value
        \State $offset \gets 0$
        \State $L \gets$ payload length
        \While{$offset < L$}
            \State Determine $s_{frag}$
            \State Align to 8-byte boundary (if non-final)
            \State Set IP.id $\gets ID$
            \State Set IP.frag $\gets offset/8$
            \State Set MF flag appropriately
            \State Append fragment to $P'$
            \State $offset \gets offset + s_{frag}$
        \EndWhile
    \Else
        \State Append $p$ to $P'$
    \EndIf
\EndFor
\State \Return $P'$
\end{algorithmic}
\end{algorithm}


During reconstruction, each fragment inherits routing-relevant metadata from the original packet. The Identification field is consistently assigned across fragments derived from the same parent packet, the Fragment Offset is computed based on cumulative payload length, and the More Fragments flag is set appropriately. Header checksums are recalculated to ensure strict RFC~791 compliance.

To maintain temporal realism, the first fragment retains the original timestamp, while subsequent fragments are assigned minimal incremental delays that emulate realistic transmission bursts without significantly altering overall flow timing characteristics.

\section{Protocol Compliance Constraints}
Strict adherence to IPv4 fragmentation rules is required to ensure that generated traffic traverses intermediate routers, firewalls, and operating system stacks without being discarded.

\subsubsection{8-Byte Alignment}

The IPv4 \texttt{Fragment Offset} field encodes offsets in units of 8 bytes. Consequently, all fragment payloads, except the final one, must be multiples of 8 bytes.

To ensure this, the algorithm enforces:

\[
s_{frag} \leftarrow \left\lfloor \frac{s_{frag}}{8} \right\rfloor \times 8
\]

for intermediate fragments.

\subsubsection{Header Field Management}

For each fragment $f_j$ derived from parent packet $p$:

\begin{itemize}
    \item \textbf{Identification (IP.id):}  
    A unique 16-bit identifier
    \[
    ID \sim \mathcal{U}[1, 65535]
    \]
    is generated per parent packet and assigned to all its fragments.

    \item \textbf{Flags:}  
    The \texttt{MF} (More Fragments) bit is set to 1 for all but the last fragment.

    \item \textbf{Fragment Offset:}  
    Computed as cumulative payload bytes divided by 8:
    \[
    \texttt{frag} = \frac{\text{payload\_offset}}{8}.
    \]

    \item \textbf{Metadata Preservation:}  
    Fields such as source address, destination address, protocol identifier, TTL, and TOS are copied unchanged.
\end{itemize}


\subsection{Stochastic Properties and Adversarial Implications}

The effectiveness of this mechanism derives from its non-deterministic nature. Identical executions on the same input trace produce distinct traffic realizations.

Three primary entropy sources are introduced:

\begin{enumerate}
    \item \textbf{Selection Entropy:}  
    A packet may be fragmented in one execution but not in another, governed by $\rho$.

    \item \textbf{Structural Entropy:}  
    The number and sizes of fragments vary due to uniform random sampling.

    \item \textbf{Identifier Entropy:}  
    Randomized IP identification values modify header checksums and signature-level patterns.
\end{enumerate}

These stochastic variations induce measurable distribution shifts in features such as packet count per flow, mean and variance of packet size, burst density and fixed-length tensor zero-padding ratios.

As a result, machine learning-based NIDS must generalize to a substantially expanded traffic manifold. Models overfitting to non-fragmented statistical patterns are particularly susceptible to misclassification under this transformation.

\subsection{Impact on NIDS Evasion}

The introduction of controlled fragmentation enhances evasion effectiveness by inducing structural distribution shift while preserving semantic equivalence at higher protocol layers.

At the packet and flow level, fragmentation modifies size distributions, increases packet counts per flow, and alters burst density patterns. These transformations directly affect statistical descriptors commonly exploited by flow-based ML-NIDS, including mean packet size, variance, and inter-arrival dynamics.

In addition, fragmentation perturbs feature extraction pipelines that rely on fixed-length windows or tensor-based representations. By increasing packet granularity and redistributing payload across multiple fragments, the method modifies zero-padding ratios, temporal ordering patterns, and aggregate flow statistics. Such perturbations can shift samples beyond previously learned decision boundaries without altering the underlying malicious functionality.

As a result, the adversarial manipulator operates not only in feature space but also in structural space. This dual-layer transformation expands the reachable adversarial manifold and exposes overfitting tendencies in ML-based NIDS models, thereby supporting more robust adversarial training and systematic evaluation.

\section{Preprocessing Pipeline}
\label{sec:preprocessing}

A robust preprocessing pipeline is essential for preparing network traffic data for adversarial manipulation and subsequent evaluation. In this thesis, preprocessing is performed using the \texttt{lucid\_dataset\_parser.py} script from the Flatten Lucid framework, which transforms raw packet capture (PCAP) files into structured datasets suitable for machine learning and adversarial analysis.

\textbf{Array Dimensions:} For the flattened model, each flow is represented as a 1D array of 21 features. For matrix input, each flow is a 2D array of shape (max\_flow\_len, num\_features), e.g., (100, 21) for 100 packets per flow and 21 features per packet. The HDF5 files store train/val/test datasets as arrays of shape (num\_flows, max\_flow\_len, num\_features) for matrix input, or (num\_flows, 21) for flattened input.

The preprocessing flow consists of the following steps:
\begin{enumerate}
    \item \textbf{Feature Extraction:} Each packet is parsed to extract a comprehensive set of features, including timestamps, packet lengths, IP fragmentation flags, protocol identifiers, and transport-layer flags.
    \item \textbf{Flow Construction:} Packets are grouped into flows based on 5-tuple identifiers and time windows.
    \item \textbf{Labeling:} Flows are labeled as benign or attack using dataset-specific rules and encoding.
    \item \textbf{Balancing and Splitting:} The dataset is balanced and split into training, validation, and test sets.
    \item \textbf{Normalization and Flattening:} Feature vectors are normalized and optionally flattened to 1D arrays.
    \item \textbf{Export to HDF5:} The processed data is exported as HDF5 files, with separate files for training, validation, and test sets.
\end{enumerate}

\textbf{Required Input for TrafficManipulator:} The TrafficManipulator framework requires as input the HDF5 datasets generated by this preprocessing pipeline. Additionally, a target benign feature set to mimic must be prepared and saved as a \texttt{.npy} file. This file should contain feature arrays classified as benign by the targeted system.

\textbf{Parameter Tuning for TrafficManipulator:} The algorithm exposes nine key parameters, divided into three groups:
\begin{itemize}
    \item \textbf{PSO velocity update:} \texttt{w} (inertia), \texttt{c1} (cognitive force), \texttt{c2} (social force)
    \item \textbf{PSO search configuration:} \texttt{max\_iter} (iterations), \texttt{particle\_num} (population), \texttt{grp\_size} (neighborhood size)
    \item \textbf{Manipulator parameters:} \texttt{grp\_size} (packets mutated per processing), \texttt{min\_time\_extend}/\texttt{max\_time\_extend} (max interarrival time extension), \texttt{max\_cft\_pkt} (max crafted packets per original), \texttt{max\_crafted\_pkt\_prob} (probability of crafting)
\end{itemize}
Example parameter settings:
\begin{verbatim}
m.change_particle_params(w=0.6, c1=0.7, c2=1.4)
m.change_pso_params(max_iter=5, particle_num=10, grp_size=5)
m.change_manipulator_params(grp_size=5, min_time_extend=0., max_time_extend=5., max_cft_pkt=4, max_crafted_pkt_prob=0.3)
\end{verbatim}

This preprocessing pipeline ensures that all subsequent adversarial experiments are conducted on standardized, well-structured datasets, enabling reproducible and meaningful evaluation of evasion techniques.

